\clearpage\begingroup\SourceLeaf{28}\NoSourceHead
\begin{center}
[B.]\\[.8em]
[MAASS UND MESSEN.]\\[1.3em]
[1.]\\[.7em]
\textsc{Gauss} an \textsc{Olbers}. [Göttingen, 8. Dezember 1817\textsuperscript{*})].
\par\medskip\rule{.78\textwidth}{.35pt}\par
{\small [\textsc{Wilhelm Olbers}, Sein Leben und seine Werke, II, 1, 1900, S. 674, 675.]}\par
\rule{.78\textwidth}{.35pt}
\end{center}
\begingroup\renewcommand{\thefootnote}{}\footnotetext{[\textsuperscript{*}) Datum des Poststempels.]}\endgroup
\noindent .\ .\ .\ .\ .\ .\ .\ .\ .\ .\ .\ .\ .\ .\ .\ .\ .\ .\ .\ .\ .
\par\noindent\hspace*{\parindent}Gleichfalls sehr interessant ist mir die Aussicht einer vielleicht allgemeinen
Einführung des französischen Maasssystems. Höchst bequem finde
ich dieses System und ich bediene mich desselben gern überall, und glaube,
dass alles oder das Meiste, was man gegen allgemeine Einführung gesagt hat,
auf Vorurtheilen beruhet. Nur bei den allerfeinsten Messungen, glaube ich,
entstehen grosse Inconvenienzen aus der Einführung eines \textls[120]{natürlichen} Maasssystems,
und man muss daneben immer irgend ein Maassindividuum haben.
Wenn z. B. im Jahr 1900 in irgend einem Welttheile eine neue Gradmessung
ausgeführt wird, so darf man den Werth des Grades nicht nach Meters angeben,
weil jede Gradmessung direct oder indirect den Zweck hat, den Meter
zu suchen; giebt man ihn nach Meters an, so bedeutet der Meter nicht
\(\frac{1}{10000000}\) Erdquadrant, sondern die Länge desjenigen Stückes Eisen, bei der
und der Temperatur, welche man im Jahr NN dafür hielt, und an dem Ort
NN in dieser Meinung deponirte. Es ist also ein nie aufhörendes Schwanken.
\par\noindent .\ .\ .\ .\ .\ .\ .\ .\ .\ .\ .\ .\ .\ .\ .\ .\ .\ .\ .\ .\ .
\clearpage\endgroup
